\section{Methodology}
\label{sec:methodology}

\subsection{Overview}
\label{sec:method-overview}

We study whether, and how, retrieval-augmented generation (RAG) improves an
LLM's ability to detect malicious PyPI packages, by comparing seven
alternative pipelines~E0--E6 that differ only in the \emph{retrieval strategy}
they apply before the final classification. At a high level, every pipeline
(1)~takes the target package's \texttt{setup.py} source as input,
(2)~builds a query representation from it,
(3)~optionally retrieves reference examples from a knowledge base of
previously seen packages,
(4)~feeds the target together with the retrieved evidence (if any) to an
instruction-tuned LLM, and
(5)~asks the model for a structured verdict: malicious or benign, a confidence
score, security-relevant behaviors with source-line citations, and a rationale.
Figure~\ref{fig:pipeline} illustrates the pipeline.

\begin{figure}[!ht]
    \centering
    \begin{tikzpicture}[
        node distance=6mm,
        box/.style={rectangle, rounded corners=2pt, draw=black!60, fill=gray!10,
                    minimum height=8mm, align=center, font=\scriptsize,
                    inner sep=3pt},
        arr/.style={-{Stealth[length=2.5mm]}, thick, black!70},
    ]
    \node[box] (target) {Target package \\ \texttt{setup.py}};
    \node[box, right=of target] (num) {Line-numbered \\ source};
    \node[box, below=of num] (beh) {AST behavior \\ summary (E4/E6)};
    \node[box, right=of num] (ret) {Retriever \\ dense / BM25 / hybrid / contrastive};
    \node[box, above=of ret] (kb) {Knowledge base \\ mal. + ben. train examples};
    \node[box, right=of ret] (prompt) {Prompt builder \\ evidence + target};
    \node[box, right=of prompt] (llm) {LLM \\ (Llama-3.1-8B)};
    \node[box, right=of llm] (out) {Structured verdict \\ \{label, conf, behaviors\}};
    \draw[arr] (target) -- (num);
    \draw[arr] (num) -- (ret);
    \draw[arr] (num) -- (beh);
    \draw[arr] (beh) -- (ret);
    \draw[arr] (kb) -- (ret);
    \draw[arr] (ret) -- (prompt);
    \draw[arr] (num) to[out=25, in=155] (prompt);
    \draw[arr] (prompt) -- (llm);
    \draw[arr] (llm) -- (out);
    \end{tikzpicture}
    \caption{Overview of the evaluated pipelines. All stages are shared across
    methods except query construction and retrieval (E4/E6 additionally build an
    AST behavior summary as the query).}
    \label{fig:pipeline}
    \Description{Pipeline with target package input, line numbering,
    optional AST behavior summary, retrieval from a knowledge base, prompt
    construction, an LLM call, and a structured verdict output.}
\end{figure}

\subsection{Dataset}
\label{sec:dataset}

We reuse the dataset assembled in the EASE~2025 study on RAG for malicious
code detection~\cite{ease2025rag}, which itself mines packages from Datadog's
malicious-software-packages dataset~\cite{datadog-dataset}. The replication
package ships pre-split JSON files: a training partition of 1{,}158 malicious
and 3{,}005 benign packages, and a test partition of 276 malicious and 753
benign packages (1{,}029 in total). Each entry contains
the package name, file list, and the raw \texttt{setup.py} content; test
entries additionally carry an AST-derived textual behavior description.
Nine malicious test entries whose \texttt{setup.py} could not be extracted
(content string \texttt{"Not Available"}) are excluded from evaluation,
leaving 1{,}020 test samples. The training partition serves as the knowledge
base; the test partition is never used for retrieval.

The malicious class is dominated by typosquatted or squatted packages
(e.g.,~\texttt{colorama} clones, \texttt{requests}-lookalikes) whose
installation code frequently hides its malicious behavior in the
\texttt{setup.py} entry point. This makes the task representative of
real-world supply-chain attacks that LLM-based tooling is expected to catch.

\subsection{Retrieval Strategies}
\label{sec:retrieval}

All methods share the same knowledge base: the raw \texttt{setup.py} sources
of the 4{,}163 training packages, indexed in separate malicious and benign
pools. We fix the number of retrieved documents to $k=3$, following the
reference design. The seven methods are:

\begin{itemize}
    \item \textbf{E0 --- No RAG.} The target source is classified directly,
    without any retrieved context. This is the controlled baseline.
    \item \textbf{E1 --- Dense RAG.} The query is embedded with a sentence
    embedding model and matched by cosine similarity against the embedded
    training documents. We use the \texttt{Qwen3-Embedding-4B} model~\cite{qwen3emb}
    and a FAISS \texttt{IndexFlatIP} store over L2-normalized vectors~\cite{vllm2023}.
    \item \textbf{E2 --- BM25 RAG.} Sparse lexical retrieval with the
    Okapi BM25 scoring function~\cite{bm252009}; no learned components.
    \item \textbf{E3 --- Hybrid RAG.} Dense (E1) and sparse (E2) rankings are
    fused with Reciprocal Rank Fusion~\cite{rrf2009},
    $\sum_r 1/(60 + \mathrm{rank}_r(d))$; no fusion weights are trained.
    \item \textbf{E4 --- Behavior RAG.} The target is represented by a static
    behavior summary---imports, suspicious API calls, and behaviors such as
    shell execution, network access, dynamic execution, or base64 decoding,
    extracted with Python's \texttt{ast} module---and dense retrieval runs over
    behavior summaries of the malicious pool, analogous to the
    AST descriptions used in the original dataset.
    \item \textbf{E5 --- Contrastive RAG.} Dense retrieval is run twice:
    top-$k$ documents from the malicious pool \emph{and} top-$k$ from the
    benign pool, and both groups are shown to the model with explicit
    \texttt{MAL}/\texttt{BEN} markers.
    \item \textbf{E6 --- Behavior + Contrastive RAG.} Behavior-summary
    queries (as in E4) combined with dual-pool contrastive evidence (as in E5).
\end{itemize}

\subsection{LLM Configuration and Prompting}
\label{sec:llm-config}

Classification is performed by \texttt{Llama-3.1-8B-Instruct}~\cite{llama3}
served through an OpenAI-compatible vLLM endpoint~\cite{vllm2023}. Generation
uses temperature $0$ and a maximum of 1{,}024 output tokens. The system prompt
instructs the model that it is performing static source-code analysis, not to
assume behavior unsupported by the source or retrieved evidence, and that
retrieved examples are references rather than proof; every claimed behavior
must cite target source lines. The target source is line-numbered before it is
embedded in the prompt. The output schema is enforced as JSON and requires
\texttt{verdict} (malicious/benign), \texttt{confidence},
\texttt{behaviors} (a list of \{type, evidence\} objects) and
\texttt{rationale}. Retrieved documents are labeled with stable identifiers
(\texttt{MAL\_*} / \texttt{BEN\_*}) and their retrieval scores.

\subsection{Evaluation Metrics}
\label{sec:metrics}

\paragraph{Classification.} We report precision, recall, and F1 for the
malicious class, benign-class F1, macro-F1, balanced accuracy, false-positive
rate (FPR), and false-negative rate (FNR). Because JSON decoding can fail
(e.g., output truncation by the serving backend), we additionally report
\emph{coverage}: the fraction of samples whose response parsed successfully;
unparsed samples are counted as failures and excluded from metric
denominators, so that success and failure modes are reported separately.

\paragraph{Explanation faithfulness.} As an approximation of whether the
model ``hallucinates'' security-relevant behaviors, we ground each claimed
behavior in a lightweight static checker built on Python's \texttt{ast}
module, which flags shell execution, process creation, network access, file
writes/deletion, dynamic execution, environment access, base64 decoding,
remote download, and persistence. Free-form behavior names produced by the
model are keyword-normalized to this vocabulary; a claim is
\emph{supported} when the corresponding static flag is present on any line of
the target. We report the unsupported-claim rate
$= \mathrm{unsupported} / \mathrm{claimed}$, behavior precision and behavior
recall. Because static analysis is neither sound nor complete on obfuscated
code, we treat these figures as an \emph{automated static-grounding proxy},
not as hallucination ground truth.

\subsection{Implementation and Reproducibility}
\label{sec:implementation}

The pipeline is implemented in Python 3.12 as a modular package
(retrieval, static analysis, prompting, evaluation modules) without
framework dependencies; BM25 uses \texttt{rank-bm25}, dense indexes use
\texttt{faiss}, and LLM/embedding calls go through a single
OpenAI-compatible client so that any endpoint can be swapped by environment
variables. All LLM responses, embeddings, and retrieval results are cached
(SQLite) keyed by prompt hash, sample content, evidence IDs, and generation
parameters; retrieval indexes are persisted on disk; and every run records
its configuration, seed, model IDs, and git commit. Run scripts for all seven
methods, the evaluation code, and raw per-sample outputs are included in the
replication package.
\endinput
